% W5i89_HardProblems.tex
% DFD 20-Agent Research Programme --- Iteration 89
% Hard Problems, Honest Negatives, and Open Questions

\documentclass[11pt,a4paper]{article}
\usepackage[utf8]{inputenc}
\usepackage{amsmath,amssymb}
\usepackage{booktabs}
\usepackage{geometry}
\usepackage{xcolor}
\usepackage{tcolorbox}
\usepackage{enumitem}
\geometry{margin=2.5cm}

\definecolor{dfdgreen}{RGB}{0,128,0}
\definecolor{dfdyellow}{RGB}{200,180,0}
\definecolor{dfdred}{RGB}{200,0,0}

\title{W5i89: Hard Problems and Honest Negatives\\[6pt]
\large PRL Referee Risk $\cdot$ N-body Systematic Floor $\cdot$ Pre-i90 Assessment Preparation}
\author{DFD 20-Agent Research Programme}
\date{2026-03-24}

\begin{document}
\maketitle

\section{PRL Referee Scenario Planning}

With 2 referees assigned, the PRL $\alpha$ letter will receive reports at i91--i93. Scenario planning:

\begin{tcolorbox}[colback=green!5!white, colframe=dfdgreen, title=\textbf{Scenario A: Both Referees Positive (15\%)}]
Accept with minor revisions. Submit revised version within 1 week. Publication by i95.
\end{tcolorbox}

\begin{tcolorbox}[colback=yellow!5!white, colframe=dfdyellow, title=\textbf{Scenario B: Split Decision (30\%)}]
One positive, one negative. Editor likely requests major revision or third referee. Response time: 2--4 weeks. If major revision: resubmit with additional material. If third referee: wait additional 4--6 weeks.
\end{tcolorbox}

\begin{tcolorbox}[colback=yellow!5!white, colframe=dfdyellow, title=\textbf{Scenario C: Both Negative, Constructive (35\%)}]
Rejection with constructive feedback. Revise and resubmit to PRD Rapid Communications within 2 weeks. The CQG acceptances provide credibility for PRD submission.
\end{tcolorbox}

\begin{tcolorbox}[colback=red!5!white, colframe=dfdred, title=\textbf{Scenario D: Both Negative, Dismissive (20\%)}]
Rejection without constructive feedback (``numerology,'' ``not physics''). Post extended version on arXiv. Submit to PRD (full paper, not Rapid). Consider JHEP or Annals of Physics as alternatives. The community response to the arXiv posting will determine the next step.
\end{tcolorbox}

\textbf{Strategic position}: Even in Scenario D, the programme is not derailed. The $c_T$ paper is accepted, BD is in revision, and 10 other papers are under review. A PRL rejection does not invalidate the theory.

\section{N-Body Systematic Floor}

With the suite complete, the systematic floor of the DFD N-body predictions can be quantified:

\begin{center}
\begin{tabular}{lcc}
\toprule
\textbf{Source of Uncertainty} & \textbf{Magnitude} & \textbf{Reducible?} \\
\midrule
Statistical (finite box, shot noise) & $0.3$--$0.5\%$ & Yes (more runs) \\
Baryonic physics model & $0.2$--$0.5\%$ & Partially (better subgrid) \\
$\psi$-field solver accuracy & $< 0.01\%$ & Yes (but already negligible) \\
Halo finder algorithm & $0.1$--$0.3\%$ & Partially (multi-finder) \\
HOD model uncertainty & $0.3$--$0.5\%$ & Partially (SAM alternative) \\
\midrule
\textbf{Total systematic floor} & $\sim 0.5$--$1.0\%$ & \\
\bottomrule
\end{tabular}
\end{center}

The DFD signal ($1.5$--$3\%$) is above the systematic floor ($0.5$--$1\%$) by a factor of 2--3. This is sufficient for Euclid detection but marginal. Future work should focus on reducing the baryonic and HOD uncertainties.

\section{Updated Hard Problems}

\begin{enumerate}
\item \textbf{PRL referee reports (CRITICAL PATH).} Expected i91--i93. All scenarios planned.

\item \textbf{N-body Paper I completion.} At 90\%. Target: 100\% at i90. Last 3 sections being drafted.

\item \textbf{Euclid pre-registration arXiv posting.} Target: i90. Paper at 78\%. Remaining work is polish, not content.

\item \textbf{$\alpha$ running discrepancy.} Unchanged. Deferred to i93+.

\item \textbf{Euclid DR1 timing.} ESA: October 2026 confirmed. No delays.

\item \textbf{JOSS review.} Software paper under review. JOSS reviews focus on code quality and documentation, not physics. Low risk.
\end{enumerate}

\section{Honest Negatives: Status at i89}

\begin{center}
\begin{tabular}{clcc}
\toprule
\textbf{\#} & \textbf{Honest Negative} & \textbf{Status} & \textbf{Change} \\
\midrule
1 & $\alpha$ residual 0.55 ppm & Active & --- \\
2 & $m_s$ tension $2.6\sigma$ & Active & --- \\
3 & $\Sigma m_\nu$ margin 2.2 meV & Active & --- \\
4 & $S_8$ partial resolution & Active & --- \\
5 & $E_G$ $2.9\sigma$ standalone & Active & --- \\
6 & Bispectrum $2.3\sigma$ & Active & --- \\
7 & 5 unfalsifiable predictions & Active & --- \\
8 & N-body $k < 0.02$ & \textbf{Resolved} & Closed at i86 \\
9 & Single-author & Active & --- \\
10 & PRL rejection $\sim 75\%$ & Active & Awaiting reports \\
11 & $w_a$ if confirmed & Active & --- \\
12 & IH weakens $\Sigma m_\nu$ & Active & --- \\
13 & Hubble tension not resolved & Active & --- \\
14 & Multi-probe $< 5\sigma$ & Active & Improving (new probes) \\
\bottomrule
\end{tabular}
\end{center}

No new honest negatives. The N-body completion has not revealed any surprises. The systematic floor is within expectations.

\section{Pre-i90 Assessment: What Must i90 Accomplish?}

The i90 milestone (90\% mark) should deliver:

\begin{enumerate}[nosep]
\item N-body Paper I at 100\% (15th paper)
\item Euclid pre-registration on arXiv
\item Mock catalogue data release on Zenodo
\item 90\% Mark Assessment: comprehensive review of what the final 10 iterations must accomplish
\item BD letter: accepted (expected)
\item CQG $c_T$: published (proofs returned)
\end{enumerate}

\end{document}
